\chapter{研究设计}

\section{访谈提纲设计}

根据前文的文献综述，国际学术流动对博士生的影响涵盖学术网络拓展、科研能力提升、隐性知识积累、学术声誉积累以及职业发展等多个维度。然而，现有研究较少对流动过程中各类学术资源的积累机制、转化路径及其对学术职业发展的影响进行深入探索。因此，本研究将通过深度访谈，利用扎根理论来对此议题进行分析。基于上述研究需求和现有相关理论开展访谈提纲设计。

通过3轮的小组讨论，形成初步访谈大纲，接着选取参与过国家留学基金委资助的国际交流项目的科研人员5人进行预访谈。对访谈结果进行汇总分析，并征求领域专家意见，修订初始访谈大纲，形成最终的《国际学术流动资本积累成效研究》访谈提纲。最终访谈提纲包括两部分：第一部分为受访者基本信息，包括研究方向、学术背景等，具体问题见表\ref{tab:basic-questions}；第二部分为访谈主体内容，围绕导师及其团队、国际科研合作网络、知识获取与转移、知识生产与创新、职业路径与发展等主题展开，访谈提纲设计以开放性问题为主导，避免引导性提问；问题编排遵循由宽泛到具体、由浅入深的原则，引导受访者自然连贯地对其国际流动经历与学术发展过程进行回溯，部分具体问题如表\ref{tab:interview-outline}所示。

\begin{table}[htbp]
\centering
\caption{受访者基本情况}
\label{tab:basic-questions}
\begin{tabular}{cl}
\toprule
序号 & 问题 \\
\midrule
1 & 您的专业领域以及您的主要研究方向是什么？ \\
2 & 您留学的开始和结束年份分别是哪一年？留学的总时长是多久？ \\
3 & 您留学的学校及学院名称是什么？ \\
4 & 您在留学期间所学习的专业是什么？ \\
5 & 您目前的工作单位与职称是什么？ \\
6 & 您是在博士几年级出去的？您有因为留学而延期毕业吗？ \\
\bottomrule
\end{tabular}
\end{table}

\begin{table}[htbp]
\centering
\caption{国际流动对博士研究生的多维度影响访谈提纲（部分）}
\label{tab:interview-outline}
\begin{tabular}{p{3cm}p{10cm}}
\toprule
主题 & 具体问题 \\
\midrule
导师及其团队 & 1. 您导师的主要研究方向或专业领域是什么 \newline 2. 导师团队有什么特点，对您是否产生影响？ \newline 4. 国外导师在哪些方面给予了您指导？指导强度是怎样的？ \\
科研合作 & 9. 在留学期间，您是否以及如何与国外导师及其团队成员进行科研合作？ \newline 10. 合作研究和产出对您个人的学术成长有着怎样的影响？ \\
知识获取与转移 & 13. 您在留学过程中获取了哪些类型的知识？如何获得？ \newline 14. 工作后，您是否将国外所学引入国内？ \\
知识生产与创新 & 15. 您的留学经历是否以及如何对您的知识生产产生了影响？ \newline 16. 您的留学经历是否促进了您的创新性产出？如何促进？ \\
职业路径与发展 & 17. 您觉得您的国际流动经历对您获得初职是否有积极作用？ \newline 18. 您觉得您的国际流动经历对您职称晋升是否有积极作用？ \\
开放性讨论 & 19. 除上述内容外，是否有其他与留学经历、学术成长、职业发展相关的信息，希望补充说明？ \\
\bottomrule
\end{tabular}
\end{table}

\section{样本选择}

扎根理论的抽样要求围绕研究问题，挑选能最大程度覆盖研究信息的对象来访谈，深度访谈需要获取丰富、深入的访谈资料，更看重访谈的质量，因此一般采用非随机抽样方式。本研究以"国际学术流动对科研人员早期学术资本积累与职业发展的影响"为主题，为确保研究对象与研究主题的契合度与研究结论的有效性，被访者的选取需满足：（1）具有六个月以上的国家留学基金委资助的留学经历；（2）目前已回国，并在国内高校、科研院所或其他学术机构从事教学、科研相关工作；（3）愿意分享其国际流动过程中的真实经历、感知与反思。通过以上标准，旨在确保访谈对象切实具有丰富的国际学术流动体验，能够为理解国际流动对学者学术生涯与个人发展的多重影响提供充分、有深度的信息。

在访谈对象的选取方面，本研究采用目的性抽样与滚雪球抽样相结合的策略，以最大限度覆盖不同学科背景、海外经历时长和回国年限的归国学者群体。其中一部分受访者通过学术关系网络由熟人推荐引入；另一部分则是从前期问卷调查中自愿接受后续访谈的研究人员中选取，并通过发送访谈邀请邮件的方式完成招募，从中依据理论饱和原则及样本多样性要求进一步选取合适对象，具体从以下角度考量：一是学科多样性，涵盖文学、理学、医学、社会学等各类学科，确保样本学科分布均衡；二是出国时间多样性，选取的出国时间覆盖2000年代、2010年代及2020年代等不同阶段，以呈现不同时期国际流动背景下的群体特征差异；三是出国国家多样性，选取前往不同国家的学者，体现不同国家环境对学者的潜在影响；四是受访者性别维度，兼顾不同性别受访者的比例与结构，以避免样本在性别构成上的单一化倾向。所有潜在受访者均通过学术邮箱或微信等方式联系，招募过程中明确告知本研究的目的、研究内容，同时说明访谈将进行录音，且访谈内容仅用于学术研究，不作其他用途，相关个人信息均予以严格保密，受访者若对任一问题不愿作答，可随时提出。受访者在知悉并同意相关内容的前提下自愿参与访谈。本研究共招募并完成22位归国学者的深度访谈，受访者的基本信息（如性别、学科领域、海外国别、留学时长等）详见表\ref{tab:interviewees}。

\begin{longtable}{cccccc}
\caption{受访者基本情况} \label{tab:interviewees} \\
\toprule
编号 & 性别 & 学科领域 & 海外国别 & 留学年份区间 & 留学时长 \\
\midrule
\endfirsthead
\caption[]{受访者基本情况（续）} \\
\toprule
编号 & 性别 & 学科领域 & 海外国别 & 留学年份区间 & 留学时长 \\
\midrule
\endhead
\bottomrule
\endfoot
A01 & 女 & 化学 & 德国 & 2013--2014 & 12个月 \\
A02 & 女 & 图书馆学 & 美国 & 2017--2018 & 12个月 \\
A03 & 女 & 化学 & 西班牙 & 2015--2017 & 14个月 \\
A04 & 男 & 计算机 & 美国 & 2017--2018 & 12个月 \\
A05 & 男 & 管理科学与工程 & 新西兰 & 2018--2019 & 12个月 \\
A06 & 男 & 计算数学 & 美国 & 2014--2016 & 22个月 \\
A07 & 男 & 化学 & 西班牙 & 2015--2016 & 12个月 \\
A08 & 男 & 法学 & 德国 & 2013--2014 & 12个月 \\
A09 & 男 & 情报学 & 美国 & 2010--2011 & 15个月 \\
A10 & 女 & 档案学 & 美国 & 2016--2017 & 12个月 \\
A11 & 男 & 人文地理学 & 美国 & 2019--2020 & 12个月 \\
A12 & 女 & 环境科学 & 澳大利亚 & 2016--2018 & 17个月 \\
A13 & 男 & 计算机 & 英国 & 2012--2012 & 6个月 \\
A14 & 女 & 生物信息学 & 美国 & 2007--2009 & 24个月 \\
A15 & 女 & 信息资源管理 & 加拿大 & 2021--2025 & 54个月 \\
A16 & 女 & 行政管理 & 美国 & 2019--2020 & 12个月 \\
A17 & 男 & 计算数学 & 美国 & 2011--2013 & 24个月 \\
A18 & 女 & 计算机应用技术 & 法国 & 2006--2007 & 12个月 \\
A19 & 男 & 戏剧影视学 & 美国 & 2014--2015 & 12个月 \\
A20 & 男 & 哲学 & 美国 & 2017--2018 & 12个月 \\
A21 & 女 & 新闻传播学 & 法国 & 2018--2019 & 12个月 \\
A22 & 男 & 遗传学 & 瑞典 & 2009--2011 & 24个月 \\
\end{longtable}

\section{数据收集}

本研究以访谈录音为主要数据收集方式，辅以访谈过程中的现场笔记，数据收集分为两个阶段：第一阶段是访谈前期准备，在受访者的个人主页、所在机构的官方网站、博士论文致谢、泛研网等数据库搜集个人资料，以获得基础信息，为后续访谈的顺利开展做好准备，目的是节省提问时间，增加后续访谈的针对性；第二阶段是正式访谈，根据访谈提纲灵活调整问题顺序，针对受访者的回答适时进行追问与延伸探讨。

因大部分受访者身处异地，本研究采用两种访谈形式。对于有条件的受访者采用面对面访谈的形式，研究者提前与受访者商定访谈时间与地点，选择安静、无外界打扰的场所，同时做好现场录制工作，保障访谈顺利进行；对于距离较远的受访者进行线上访谈，研究者提前与受访者沟通并征得同意后，以语音或视频形式开展，提前确认访谈时间、线上沟通软件，确保网络通畅，且对访谈过程进行云端录制，以完整保留访谈内容用于后期分析。

考虑到访谈时长会直接影响数据的丰富度，时长过短易导致研究数据收集不充分，为保障受访者能够充分回顾、完整表达相关内容，结合访谈提纲及预访谈的结果与反馈，本研究将单次访谈时长设定为60分钟左右，并在访谈正式开始前明确告知受访者访谈的大致时长。此外，在访谈之前，研究者告知受访者访谈目的和访谈数据保密措施，以确保受访者能提供真实信息。在访谈过程中，研究者需把握访谈的时间、话题与节奏，积极引导受访者主动叙述、充分表达，避免话题长时间偏离主题；同时保持中立立场，不进行主观引导，不表露个人态度与观点。若受访者对研究问题存在理解偏差或提出疑问，访谈者需及时对问题进行阐释，灵活调整表述语气与措辞，确保双方对问题的理解达成一致；当挖掘到具有研究价值的信息时，访谈者需进一步追问探究，引导受访者展开详述，以获取更具深度、可供后续深入研究的话题与内容。访谈结束后，首先向受访者表示感谢，同时告知其可随时与研究团队联系补充相关内容，也说明研究团队若在内容分析中需进一步核实问题，可能会再次与对方沟通；对有条件的受访者，同步邀请其推荐其他符合研究要求的潜在受访对象。

访谈工作于2024年6月启动，2026年3月完成，累计形成26.2小时访谈语音资料，单次访谈平均时长为1.2小时\cite{sunxiaoe2011}。访谈结束后，对全部语音资料开展转录与处理，最终形成约23万字的访谈文字资料。同时，每次访谈完成后，均对访谈过程及提纲设计进行反思与优化，并结合访谈中获取的有价值信息，为研究增设新的分析维度。

\section{数据编码}

本研究采用NVivo 20质性分析软件对访谈文本进行系统整理与编码分析，同时辅以Word和Excel记录研究笔记与整理分析内容，遵循扎根理论的分析路径，对资料进行逐级抽象与理论建构。资料分析依次经历开放式编码、主轴编码与选择性编码三个阶段，在持续比较与反复归纳的过程中逐步提炼核心范畴。

整个编码过程分为两个阶段。第一阶段为资料驱动的开放式编码，在未预设具体理论框架的情况下，对22份访谈文本进行逐句阅读与意义单元提取，尽量贴近受访者原始表达生成初始概念，并通过持续比较对相似概念进行合并与归类。第二阶段在前期归纳结果的基础上，引入理论视角进行进一步整合与提升，将既有概念与范畴进行重新梳理与结构化分析，最终形成以布迪厄资本理论为总体分析框架，结合格兰诺维特强弱关系理论与路桑斯心理资本理论进行解释深化的编码体系。

\subsection{开放式编码}

开放式编码是质性研究分析的第一步，其主要目的是对访谈资料进行逐句阅读和意义提取，将受访者的原始表述转化为具有分析意义的概念。在初始概念提炼阶段，本研究尽量保留受访者的本土话语逻辑，不对原始表述进行强制归类，而是紧贴材料本身，通过持续比较不同访谈文本中的相似表达，将内容相近的概念逐步归并，形成初步的类别结构。在具体操作过程中，本研究首先对访谈文本中的原始语句进行逐句分析，并按照"初始标签——初始代码——概念"三个层级进行整理，再对生成的原始概念进行梳理与整合。其中，初始标签是对原始语句进行最贴近材料的标注，通常采用受访者的原话或对其表述的直接概括，用以反映该语句所表达的具体行为、感受或经验；初始代码是在初始标签基础上进一步归纳与提炼的结果，通过整合表达相似意义的标签，使零散表述上升为具有一定解释力的分析单元；概念则是在对多个初始代码持续比较与整合的基础上形成的更高层次类别，用以概括具有相似性质的经验现象。

在编码过程中，由于本研究的编码体系层次较多，在呈现开放式编码结果时，需要对部分概念进行理论抽象与整合。例如，原本提取出了"教学能力培养""科研方法与技能提升""研究效能提升""语言能力提升"等概念，共同指向研究者"可操作能力"的提升维度。经反复比较与归纳，将其统一整合为"学术技能提升"概念。

此外，在第一轮编码时，本研究严格遵循"从数据中生长概念"的原则，形成了大量以受访者原话为基础的初始概念。但在后续的概念整合过程中，发现部分初始概念归属于资本理论的既有范畴，例如，"产出型客观化文化资本"的概念在理论上承接布迪厄"客观化文化资本"的经典范畴，指以物质载体形式存在的文化产品，在数据中，受访者提及的"论文发表""专利产出""实验系统搭建""实验室设备引进"等经验现象，是客观化文化资本的具体表现。在对这些现象进行归类时，研究者发现客观化文化资本呈现出两种不同的性质：一类是研究者主动创造、从无到有的成果（如论文、技术平台）；另一类是研究者获取并持有的资源（如软件、数据库），据此，本研究将前者命名为"产出型客观化文化资本"，以区别于后者的"使用型客观化文化资本"。因此，在最终的开放式编码表格中，部分概念名称采用了理论术语。通过上述过程，本研究在开放编码阶段最终从访谈资料中提炼出55个概念，形成了19个范畴。表\ref{tab:open-coding}选取来自不同受访者的代表性语句，展示了从原始语句到概念生成的具体过程。

{\small
\begin{longtable}{p{6.5cm}p{2.5cm}p{2.5cm}p{2.5cm}}
\caption{开放编码阶段原始语句与概念生成示例} \label{tab:open-coding} \\
\toprule
原始语句 & 初始标签 & 初始代码 & 概念 \\
\midrule
\endfirsthead
\caption[]{开放编码阶段原始语句与概念生成示例（续）} \\
\toprule
原始语句 & 初始标签 & 初始代码 & 概念 \\
\midrule
\endhead
\bottomrule
\endfoot
我使用过他们图书馆的一些资源，就是可以看到一些我们以前就是在国内没有办法看到的一些专业书，也可能国内太贵了。如果要购买电子版也非常的贵，但是我可以直接去图书馆进行复印，拿我的专业知识书。 & 使用图书馆资源；国内未覆盖的专业书；国内专业书价格高昂；电子版价格也高；可直接复印书籍 & 学术材料获取 & 使用型客观化文化资本 \\
我觉得从我自身的这个经历来看，关于研究主题因为它的变化是很快的，我觉得更重要的是这种研究方法以及研究基础的这种夯实，我觉得对于我的影响是比较大的。 & 从自身经历出发；研究主题变化很快；夯实研究方法、研究基础影响较大 & 科研方法与技能提升 & 学术技能提升 \\
语言方面肯定是有的。就是敢说，本来我英语应该也还行，然后去了之后在口语表达方面肯定是提升了，因为我当时租房子也是一个当地人，就刻意地让自己去多锻炼一下。 & 敢于表达；和当地人主动接触锻炼语言 & 口语表达信心提升；主动增加日常口语实践机会 & 跨文化沟通信心提升；主动性 \\
对于科研的本质，我最大的收获是一个交流会上，我至今都记得讲到的一个观点就是Research is all about communication. 所有跟研究相关的，不管是学术会议交流，还是说你发表论文，其实都是跟学术界的一种交流和对话，我觉得越来越深有感触。 & 对于科研本质的认识；在交流会上收获最大；研究是与学术界的交流与对话 & 参与学术交流科研本质认知；学术理念内化 & 学术认知拓展 \\
其实这个东西是更深远的一个影响，就是说我把这个技术引进回来了以后，我们可以一直使用，对于我们的实验有更好的效果。 & 引进技术使用延续性强；获得更好实验效果 & 科研效能提升 & 学术技能提升 \\
最大的问题就是语言的问题。因为我当时去的是塞维利亚，是西班牙南部的一个城市，会说英语的人非常少，包括我们研究所里面一些博士生的英语口语都非常差。 & 语言问题大；会说英语的人少 & 语言障碍大 & 留学国家语言环境 \\
在当时当地不管你是从事哪个方向的中国人，就那么十几个人，大家就会形成那种远远跟其他的情况不太一样的，哪怕直到现在我们经常还是互相都在关心彼此 & 留学时形成中国人之间的独特联系；至今互相关心 & 海外同胞关系 & 同胞关系 \\
现在我们再出去说我某年某月在这个老师那边待过，还是很有用的，因为这个老师他本身的学术声誉和知名度很高。你从某种意义上说，你就是从这个分支下面分出来的一个人嘛。 & 提及老师的有用性；导师学术声誉和知名度高 & 导师学术声誉向学生转化 & 导师声誉 \\
一个人坐冷板凳的能力我觉得得到了很大的提升，一个人如何在孤独的情况下开展自己的工作，你存在语言隔阂、文化隔阂，然后还有背景知识的隔阂，包括生活上。 & 提升坐冷板凳能力；孤独情况下开展工作；语言隔阂、文化隔阂、生活困难 & 应对孤独能力提升 & 孤独适应能力提升 \\
在开会的时候会遇见一些从美国高校的一些学生，也不是合作论文，当时大家都在找工作，他们比我找的早一点，然后就找他们打听一下找工作的事。 & 开会时结识美国学生 & 会议结识 & 学术会议结识 \\
我所有的这些文章都是围绕着这个课题做的，课题就是在德国选定的。 & 所有文章围绕留学时选定课题 & 研究课题延续 & 研究方向延续 \\
\end{longtable}
}

在形成概念之后，本研究进一步对开放编码阶段提炼出的概念进行整理和归类。对于含义相近、反映同一类经验或现象的概念进行合并，并在此基础上归纳形成若干范畴。本研究在反复比较原始材料的基础上，参照布迪厄的资本理论、路桑斯等人的心理资本理论以及格兰诺维特的强弱关系理论对范畴进行命名与界定。例如，心理资本相关范畴的提炼参照了路桑斯提出的自我效能感、韧性、乐观、希望四个维度，结合受访者的实际表述归并形成韧性、正确科研观构建、自我效能感等范畴；社会资本相关范畴的划分参照了格兰诺维特的强弱关系理论，将受访者的社会关系区分为强关系网络与弱关系网络两类。这一做法体现了演绎与归纳相结合的编码逻辑，既忠实于受访者的原始表述，又借助理论框架赋予范畴更清晰的概念边界。表\ref{tab:open-coding-categories}呈现了开放编码阶段形成的范畴、包含的概念以及对概念的具体阐释示例，全部概念列表见附件。

{\small
\begin{longtable}{p{2.5cm}p{3.5cm}p{8cm}}
\caption{开放编码阶段形成的范畴、概念及概念内涵示例} \label{tab:open-coding-categories} \\
\toprule
范畴 & 概念 & 概念内涵 \\
\midrule
\endfirsthead
\caption[]{开放编码阶段形成的范畴、概念及概念内涵示例（续）} \\
\toprule
范畴 & 概念 & 概念内涵 \\
\midrule
\endhead
\bottomrule
\endfoot
具身化文化资本 & 学术技能提升 & 指博士生通过参与海外科研实践与学术训练，在多个技能维度上获得实质性提升，具体包括教学能力的培养、科研方法与技能的习得、研究效能的提升以及语言能力的增强。 \\
 & 学术规范与经验内化 & 指博士生在留学期间的科研实践与学术互动中，逐渐将学术规范和研究经验转化为自身较为稳定的学术意识与研究习惯的过程，具体包括学术规范的内化、学术经验的积累以及学术品味的建立。 \\
客观化文化资本 & 产出型客观化文化资本 & 指博士生主动创造的、以物质或准物质形态存在的学术成果，包括技术平台的搭建、教学课件的制作以及论文、专著、专利等学术产出。 \\
 & 使用型客观化文化资本 & 指博士生在留学期间获取并使用的外部学术资源，包括文献资料、数据与计算资源、实验设备与科研工具等。 \\
强关系网络 & 师生关系 & 指博士生与海外导师之间在长期学术合作与日常互动中建立的稳定、紧密的关系。 \\
 & 同胞关系 & 指博士生与海外同胞学者或学生之间因共同的语言背景、文化认同与相似处境而建立的紧密的支持性关系。 \\
正确科研观构建 & 科研态度 & 指博士生在国际学术流动过程中，对科研目标、研究价值和学术职业形成的较为理性和稳定的理解，表现为更加重视研究质量、学术对话和科研工作的长期价值。 \\
自我效能感 & 学术自信心提升 & 指博士生通过海外科研训练、学术互动和阶段性成果积累，对自身学术能力、研究水平和国际学术参与能力形成更为积极的自我评价。 \\
乐观 & 积极归因 & 指博士生在面对研究困境与客观限制时，能够将研究受阻的原因归结为外部条件的缺失而非自身能力的不足，避免将客观障碍内化为对自我的否定，维护积极的科研心态。 \\
留学环境特征 & 导师特征 & 指留学期间国外导师的个人特质，如导师指导特征、博士生与导师之间的研究方向相似度及文化差异程度等。 \\
\end{longtable}
}

\subsection{主轴编码}

主轴编码是在开放编码的基础上，通过系统分析各范畴之间的逻辑关系，将分散的范畴重新归类整合，形成若干具有内在关联的主范畴的过程。本研究共梳理与归纳出"文化资本"、"社会资本"、"象征资本"、"心理资本"、"国际流动特征"、"职业发展"六个主范畴，表\ref{tab:axial-coding}展示了主轴编码阶段形成的6个主范畴、对应的范畴及其内涵。在此基础上，为更清晰地说明各主范畴的基本含义，本文将其整理为表\ref{tab:axial-coding-def}。

{\small
\begin{longtable}{p{2cm}p{2.5cm}p{9.5cm}}
\caption{主轴编码形成的主范畴、范畴及内涵} \label{tab:axial-coding} \\
\toprule
主范畴 & 范畴 & 范畴内涵 \\
\midrule
\endfirsthead
\caption[]{主轴编码形成的主范畴、范畴及内涵（续）} \\
\toprule
主范畴 & 范畴 & 范畴内涵 \\
\midrule
\endhead
\bottomrule
\endfoot
国际流动特征 & 留学环境特征 & 指博士生在国际学术流动中所处的外部环境条件，主要包括留学国家语言环境、留学机构特征以及导师特征。 \\
 & 个人特征 & 指博士生自身所具有的个体差异条件，主要包括目标明确性、主动性和学科背景。 \\
 & 流动时长 & 指博士生在海外学习和研究的持续时间。 \\
文化资本 & 具身化文化资本 & 指博士生在国际学术流动过程中，通过学习、科研训练和跨文化交流逐步内化于自身的知识、能力、思维方式、学术规范和研究经验，这类资本以个体为载体，无法与持有者分离。 \\
 & 客观化文化资本 & 指博士生在国际学术流动过程中获取、使用或产出的具有学术价值的物质性或成果性资源，包括文献资料、设备工具以及论文、专利、课件、技术平台等，这类资本不依附于特定个体，可在学术场域中自由流通、转让与共享。 \\
社会资本 & 强关系网络 & 指在留学期间建立的互动频繁、情感深厚、互惠承诺紧密的社会关系，主要包括师生关系、同门关系、同胞关系。 \\
 & 弱关系网络 & 指流动者建立的互动频率较低、情感联结较弱的社会关系，主要包括通过学术会议结识、社群活动结识、偶然结识的人际关系。 \\
象征资本 & 留学身份声望 & 指博士生凭借CSC资助留学的经历与身份标签，在国内学术场域中获得的附加声誉与认可。 \\
 & 关联性声望 & 指通过与知名导师或高水平机构产生关联而获得的间接性声望。 \\
心理资本 & 自我效能感 & 指博士生在国际学术流动过程中，通过科研实践和跨文化交流逐步形成的对自身学术能力与沟通能力的积极评价和信心。 \\
 & 正确科研观构建 & 指博士生在国际学术流动过程中，通过接触多元学术文化与科研生态，逐步形成更为成熟的科研态度、科研专注力和内在驱动力。 \\
 & 韧性 & 韧性是指博士生在面对陌生环境、学术压力和生活挑战时，保持心理稳定、积极调适并持续推进研究的能力。 \\
 & 乐观 & 指博士生在面对研究困境与客观限制时，能够将研究受阻的原因归结为外部条件的缺失而非自身能力的不足，避免将客观障碍内化为对自我的否定，维护积极的科研心态。 \\
职业发展 & 初职获得 & 指博士生完成学业后，首次进入高校或科研机构担任教职或研究职位的过程。 \\
 & 职称晋升 & 指博士生从初级职称（讲师、助理研究员）向副高级（副教授、副研究员）及正高级（教授、研究员）职称的晋升过程。 \\
 & 项目申请 & 指博士生申报国家级（如国家自然科学基金、国家社会科学基金）、省部级及其他科研项目的过程。 \\
 & 机会拓展 & 指博士生因留学经历而获得的额外职业发展机会。 \\
\end{longtable}
}

\begin{table}[htbp]
\centering
\caption{主轴编码形成的六个主范畴及其定义}
\label{tab:axial-coding-def}
{\small
\begin{tabular}{p{2.5cm}p{11.5cm}}
\toprule
主范畴 & 定义 \\
\midrule
国际流动特征 & 指博士生在国际学术流动过程中所处的情境条件及其个体基础，主要包括留学环境特征、个人特征和流动时长。 \\
文化资本 & 指博士生在国际学术流动过程中，通过学习、科研训练、学术交流和研究实践所形成的知识、能力、经验与成果。该范畴既包括内化于个体的学术认知、思维方式、研究技能和学术规范，也包括以论文、专利、课件、文献资料、数据资源和科研工具等形式呈现的外在学术资源与成果。 \\
社会资本 & 指博士生在国际学术流动过程中，通过与导师、同门、同胞及其他学者建立联系而形成的关系网络，主要体现为强关系网络与弱关系网络两类。 \\
象征资本 & 指博士生在国际学术流动过程中，因留学身份、海外导师声誉和留学机构声誉而获得的认可、声望与合法性。 \\
心理资本 & 指博士生在国际学术流动过程中形成和增强的积极心理资源，主要包括科研观念、自我效能感、乐观和韧性等方面。 \\
职业发展 & 指博士生在完成学业后，在学术职业道路上获得和拓展发展机会的过程。结合本研究，职业发展主要体现在初职获得、职称晋升、项目获得和机会拓展四个方面。 \\
\bottomrule
\end{tabular}
}
\end{table}

\subsection{选择性编码}

选择性编码旨在探讨主范畴之间的关系，并通过故事线将各变量整合为统一的理论框架。在主轴编码确立六个主范畴及其内涵的基础上，选择性编码进一步揭示了这些主范畴之间的逻辑联系，具体操作过程是如表\ref{tab:selective-coding}所示。总体来看，国际流动特征主要反映博士生开展国际学术流动时所处的条件情境，文化资本、社会资本、象征资本和心理资本是其在国际流动过程中积累的主要资本类型，职业发展则体现为这些资本进一步发挥作用后的结果。基于此，本文形成了"国际流动特征——学术资本积累——职业发展"的理论模型，具体如图\ref{fig:model}所示。

{\small
\begin{longtable}{p{5.5cm}p{3cm}p{3.5cm}p{2cm}}
\caption{代码之间的关系指向性及所属范畴举例} \label{tab:selective-coding} \\
\toprule
部分原始语句 & 代码 & 代码指向性 & 主范畴 \\
\midrule
\endfirsthead
\caption[]{代码之间的关系指向性及所属范畴举例（续）} \\
\toprule
部分原始语句 & 代码 & 代码指向性 & 主范畴 \\
\midrule
\endhead
\bottomrule
\endfoot
其中一个导师，指导主要是具体的，一开始我做核磁共振，他就会教我核磁共振的一些原理……他大概有一个多月的时间，每天都是跟我一起做实验……学到一些有机化合物提取、纯化，然后结构解析这一块，算是实验方面的技能。 & 导师具体、高频的指导；实验技能提升 & 导师具体、高频的指导$\to$实验技能提升 & 国际流动特征；文化资本 \\
语言能力刚开始也不太足以支持我去很快速地融入大家，因为法国年轻人他们的语速非常快。 & 语言障碍；不利于融入团队 & 语言障碍$\to$不利于融入团队 & 国际流动特征；社会资本 \\
当时我导师带着我改过的论文，在论文写作方面的指导，应该是有助于后面的发表的……之前我们可能写作这一块没找到那个道道，他跟我讲到了那个逻辑什么的，后面回国的时候那些论文陆续都发表了。 & 导师指导；论文发表 & 导师指导$\to$论文发表 & 社会资本；文化资本 \\
当时招聘时有要求，要有国际留学经历，原则上有，但是没有规定说年限、时长。 & 国际留学经历；初职招聘硬性要求 & 国际留学经历$\to$初职招聘硬性要求 & 象征资本；职业发展 \\
其实我的自然科学基金的青年项目当时是在澳洲受到的一个启发，就是听他们的讲座交流，他们也经常会请nature、science的那些编辑来，当时有个老师就讲到的遥感，是空间信息技术的一些数据和技术的在他们环境领域的应用，然后我当时就结合了最核心的一个观点……然后把它拓展到我这个领域来，我觉得作为一个核心的创意点拿到的自科青年。 & 讲座中的观点认知；项目核心创意点 & 讲座中的观点认知$\to$项目核心创意点 & 文化资本；职业发展 \\
\end{longtable}
}

% \begin{figure}[htbp]
% \centering
% \includegraphics[width=0.9\textwidth]{media/image1}
% \caption{国际流动对早期学术资本积累和职业发展的影响模型}
% \label{fig:model}
% \end{figure}

\subsection{理论饱和度检验}

在扎根理论的数据搜集与分析过程中，当访谈资料不能再提供新的主题、属性、关系、概念时，则认为资料已经饱和\cite{yangliping2022}。本研究编码到第18个受访者时，不再出现新的概念。为验证理论是否饱和，接着又对4位科研工作者进行访谈，经分析，未出现新的概念和范畴，各主范畴（国际流动特征、文化资本、社会资本、象征资本、心理资本、职业发展）内部未出现新的属性，范畴之间的关系也未发生实质性变化。据此可以判断，本研究的理论饱和度良好，编码框架具有较高的稳定性和解释力。
